% document class
\documentclass[journal,12pt,onecolumn,draftclsnofoot]{IEEEtran}

%%%%%%%%%%%%%%%%%%%%%%%%%%%%%%%%%%%%
% packages and formatting stuff
%%%%%%%%%%%%%%%%%%%%%%%%%%%%%%%%%%%%

%\usepackage{cite}
\usepackage{setspace}
\usepackage{subcaption}
\usepackage{xcolor}
\usepackage{amsmath}
\usepackage{amssymb} %more math
\usepackage{float} % force images to be where you actually put them
\usepackage{pdfpages} % add pdfs
\usepackage[most]{tcolorbox} % frames around images


%%%%%%%%%%%%%%%%%%%%%%
%%% inserting code %%%
%%%%%%%%%%%%%%%%%%%%%%
\usepackage{listings}
\usepackage{color}

\definecolor{mygreen}{rgb}{0,0.6,0}
\definecolor{mygray}{rgb}{0.5,0.5,0.5}
\definecolor{mymauve}{rgb}{0.58,0,0.82}
\lstset{ %
	frame=none,
	backgroundcolor=\color{white},   % choose the background color
	basicstyle=\footnotesize,        % size of fonts used for the code
	breaklines=true,                 % automatic line breaking only at whitespace
	captionpos=b,                    % sets the caption-position to bottom
	commentstyle=\color{mygreen},    % comment style
	escapechar=\%,					 % if you want to add LaTeX within your code
	keywordstyle=\color{blue},       % keyword style
	stringstyle=\color{mymauve},     % string literal style
	showstringspaces=false,
}
%%%%%%%%%%%%%%%%%%%%%%%%%%
%%% BIBLIOGRAPHY STUFF %%%
%%%%%%%%%%%%%%%%%%%%%%%%%%
\usepackage[backend=biber,style=ieee,urldate=iso,date=iso,seconds=true]{biblatex}
\bibliography{eece_resources.bib}

%%%%%%%%%%%%%%%%%%%%%%
%%% PAGE FORMATING %%%
%%%%%%%%%%%%%%%%%%%%%%
\addtolength{\oddsidemargin}{-1.cm}
\addtolength{\textwidth}{2cm}
\addtolength{\topmargin}{-2cm}
\addtolength{\textheight}{3.5cm}

\def\changemargin#1#2{\list{}{\rightmargin#2\leftmargin#1}\item[]}
\let\endchangemargin=\endlist

\usepackage[margin=1in]{geometry}

\renewcommand\IEEEkeywordsname{}
\renewcommand\abstractname{}
\newcommand{\HRule}{\rule{\linewidth}{0.5mm}}

\usepackage[margin=1in]{geometry}
\linespread{1.1} %  linespacing
\renewcommand\IEEEkeywordsname{}
\renewcommand\abstractname{}

% correct bad hyphenation here
\hyphenation{op-tical net-works semi-conduc-tor}


%%%%%%%%%%%%%%%
%%% MACROS %%%%
%%%%%%%%%%%%%%%
% adding figures (path, caption, label)
% eg. \pic{path/to/image.png}{some nice caption text}{fig:image_example}

\newcommand{\pic}[3]{
	\vspace{2.5mm}
	\begin{center}
		\tcbox[sharp corners, boxrule=0.5mm, boxsep=0mm, colframe=black, colback=white, width=0.7\textwidth]{\includegraphics[width=0.85\textwidth]{#1}}
		\captionof{figure}{#2}
		\label{#3}
	\end{center}
	\vspace{2.5mm}
	}

%images next to each other
% eg.
%	\vspace{2.5mm}
%	\begin{figure}
%		\subpic{path/to/image.png}{some nice caption text}{fig:image_example}
%		\subpic{path/to/image2.png}{some nice caption text}{fig:image2_example}
%	\end{figure}
%	\vspace{2.5mm}

\newcommand{\subpic}[3]{
	\begin{subfigure}[h]{0.5\textwidth}
		\centering
		\tcbox[sharp corners, boxrule=0.5mm, boxsep=0mm, colframe=black, colback=white, width=0.7\textwidth]{\includegraphics[width=\textwidth]{#1}}
		\caption{#2}
		\label{#3}
	\end{subfigure}
	}

%smaller oscilloscope images
% eg. \oscpic{path/to/image.png}{some nice caption text}{fig:image_example}

\newcommand{\oscpic}[3]{
	\vspace{2.5mm}
	\begin{center}
		\tcbox[sharp corners, boxrule=0.5mm, boxsep=0mm, colframe=black!, colback=white, width=0.7\textwidth]{\includegraphics[width=0.7\textwidth]{#1}}
		\captionof{figure}{#2}
		\label{#3}
	\end{center}
	\vspace{2.5mm}
	}

%inserting python
	% eg. \inputpython{path/to/code.py}{caption about the script}{code:label}
	
\newcommand{\inputpython}[3]{
	\vspace{2.5mm}
	\lstinputlisting[language=python,caption=#2,label=#3]{#1}
	\vspace{2.5mm}
	}

%%%%%%%%%%%%%%%%%%%%%%%%%%%%%%%%%%%%
% the actual document
%%%%%%%%%%%%%%%%%%%%%%%%%%%%%%%%%%%%

\begin{document}

	%%%%%%%%%%%%%%%%%%%%%%%%%%%%%%%%%%%%%%%%%%%%%%%%%%%%%
	% Coverpage and plagiarism statement (use .pdf names)
	%%%%%%%%%%%%%%%%%%%%%%%%%%%%%%%%%%%%%%%%%%%%%%%%%%%%%

	\begin{titlepage}

	\begin{center}
		%the above is justification
		% Upper part of the page
		\includegraphics[width=1\textwidth]{./up-logo.jpg}\\[0.4cm]
		\textsc{\LARGE Department of Electrical, Electronic and Computer Engineering}\\[1.5cm]
		\textsc{\Large EMS 310 }\\[0.5cm]
		% Title
		\HRule \\[0.4cm]
		{ \huge \bfseries Practical 2: FM Modulation and Demodulation}\\[0.4cm]
		\HRule \\[0.4cm]


		\begin{minipage}{0.4\textwidth}
			\begin{flushleft} \large
				\emph{Author:}\\
			\end{flushleft}
		\end{minipage}
		\begin{minipage}{0.4\textwidth}
			\begin{flushright} \large
				\emph{Student number:} \\

			\end{flushright}
		\end{minipage}


		% Student Credentials
		\begin{minipage}{0.4\textwidth}
			\begin{flushleft} \large
				%\emph{Author:}\\
				\textsc{Kevin Nel}\\ % Put the name of the  Student here
			\end{flushleft}
		\end{minipage}
		\begin{minipage}{0.4\textwidth}
			\begin{flushright} \large
				17003769 %put Student Number here
			\end{flushright}
		\end{minipage}




		% Student Credentials
		\begin{minipage}{0.4\textwidth}
			\begin{flushleft} \large
				\textsc{Tiehan Nel} % Put the name of the  Student here
			\end{flushleft}
		\end{minipage}
		\begin{minipage}{0.4\textwidth}
			\begin{flushright} \large
				17330999 %put Student Number here
			\end{flushright}
		\end{minipage}


		% Student Credentials
		\begin{minipage}{0.4\textwidth}
			\begin{flushleft} \large
				\textsc{Anro v.d. Merwe} % Put the name of the  Student here
			\end{flushleft}
		\end{minipage}
		\begin{minipage}{0.4\textwidth}
			\begin{flushright} \large
				17004812 %put Student Number here
			\end{flushright}
		\end{minipage}


		\vfill
		% Bottom of the page
		{\large \today}
	\end{center}
\end{titlepage}


	\footnotesize

	%%
%%  University of Pretoria - Declaration of Originality
%%  Version: S4726/09 (amended 2010-07-23)
%%

\newpage

\thispagestyle{empty}
{
\renewcommand{\baselinestretch}{1}
\sffamily
\begin{center}
\textbf{\Large DECLARATION OF ORIGINALITY}
\end{center}
\begin{center}
\textbf{\large UNIVERSITY OF PRETORIA}
\end{center}

The University of Pretoria places great emphasis upon integrity and
ethical conduct in the preparation of all written work submitted for
academic evaluation.

While academic staff teach you about referencing techniques and how to
avoid plagiarism, you too have a responsibility in this regard. If you
are at any stage uncertain as to what is required, you should speak to
your lecturer before any written work is submitted.

You are guilty of plagiarism if you copy something from another
author's work (e.g. a book, an article or a website) without
acknowledging the source and pass it off as your own. In effect you
are stealing something that belongs to someone else. This is not only
the case when you copy work word-for-word (verbatim), but also when
you submit someone else's work in a slightly altered form (paraphrase)
or use a line of argument without acknowledging it. You are not
allowed to use work previously produced by another student. You are
also not allowed to let anybody copy your work with the intention of
passing if off as his/her work.

Students who commit plagiarism will not be given any credit for
plagiarised work. The matter may also be referred to the Disciplinary
Committee (Students) for a ruling. Plagiarism is regarded as a serious
contravention of the University's rules and can lead to expulsion from
the University.

The declaration which follows must accompany all written work
submitted while you are a student of the University of Pretoria. No
written work will be accepted unless the declaration has been
completed and attached.

\begin{center}
\begin{tabular}{ll}
                        &                             \\
 Full names of student: & \makebox[3.5in]{\hrulefill} \\ \\
 Student number:        & \makebox[3.5in]{\hrulefill} \\ \\
 Topic of work:         & \makebox[3.5in]{\hrulefill} \\ \\
\end{tabular}
\end{center}

\textbf{Declaration}
\begin{enumerate}
\item I understand what plagiarism is and am aware of the University's
  policy in this regard.
\item I declare that this assignment report is my own original work.
  Where other people's work has been used (either from a printed source,
  Internet or any other source), this has been properly acknowledged and
  referenced in accordance with departmental requirements.
\item I have not used work previously produced by another student or any
  other person to hand in as my own.
\item I have not allowed, and will not allow, anyone to copy my work with
  the intention of passing it off as his or her own work.
\end{enumerate}

\vspace*{1cm}

SIGNATURE: \makebox[3in]{\hrulefill} \quad DATE: \makebox[1.5in]{\hrulefill}
}

\newpage



	%\includepdf[pages=-]{signatures.pdf} % add scanned in signature page

	\normalsize

	\clearpage

	\thispagestyle{empty}
	\tableofcontents

	\clearpage

	\setcounter{page}{1} %sets page counter only after toc


	\title{Phase Locked Loop}

	% authors
	%Put in Student credentials here as well
	\author{Tiehan~Nel~\IEEEmembership{12345678},
			Anro~v.d.~Merwe~\IEEEmembership{12345678},
			Kevin~Nel~\IEEEmembership{17003769}% <-this % stops a space


	\thanks{Tiehan~Nel, Anro~v.d.~Merwe, and Kevin~Nel are EMS 310 students in group 5 with the Department of Electrical, Electronic and Computer Engineering at the University of Pretoria.}}% <-this % stops a space


	% The paper headers
	\markboth{Phase Locked Loop}%
		{Shell \MakeLowercase{\textit{et al.}}: Bare Demo of IEEEtran.cls for IEEE Journals}

	\maketitle


	%%%%%%%%%%%%%%%%%%%%%%%%%%%%%%
	% abstract and keywords
	%%%%%%%%%%%%%%%%%%%%%%%%%%%%%%

	\begin{abstract}
		\begin{changemargin}{-8.55mm}{-8.55mm}
		\textbf{\textit{Abstract}}---Abstract goes here.
	\end{changemargin}
	\end{abstract}



	\begin{IEEEkeywords}
		\vspace{-2.25em}
		\begin{changemargin}{-8.55mm}{-8.55mm}
			\textbf{\textit{Keywords}}---keywords
		\end{changemargin}
	\end{IEEEkeywords}




%%%%%%%%%%%%%%%%%%%%%%%%%%%%%%
% the real stuff
%%%%%%%%%%%%%%%%%%%%%%%%%%%%%%

	% place sections here
	\section{Introduction}

	\clearpage
	\section{Aim}

	\clearpage
	\section{Theoretical Analysis}

	\clearpage
	\section{Design.tex}
	
	

	%use the macros provided inmain.tex to add images:
% eg.
% \pic{path/to/image.png}{some caption tex}{fig:figure_label} 
% 
% note:
% when referencing images in text:
% ... as seen in figure~\ref{fig:figure_label}~
%

\section{Simulations}
	
	
	

	\clearpage
	\section{Results}
% use same method of adding pictures as in Simulations.tex

	\clearpage
	\section{Discussion}


	\clearpage
	\section{Conclusion}





\printbibliography

%%%%%%%%%%%%%%%%%%%%%%%%%%%%%%
% Appendices
%%%%%%%%%%%%%%%%%%%%%%%%%%%%%%
\clearpage
\appendices
\input{Appendix.tex}

\end{document}
